\section{模型评价、局限与结论}
第一问的交互代理在训练内嵌套折及已观察的真实配方检验组上优于线性基线，但 A6--A11 曾影响模型形式判断；估算 10B/70B 表反而削弱跨规模排序主张。第二问的 B1 同源拟合精度高，B7 质量扩展经留等级检验，但半合成来源与 A/B 缺乏共同实验，质量及配比桥接只能解释为条件情景。第三问的固定策略分段凸求解、有限配方枚举与数值证书验证了模型内部优化，不能证明真实训练的无条件最佳配置。第四问同格贡献闭合与四窗口前沿回测支持条件分析；格内数据量和选择混杂、跨 Loss 坐标未知及长期外推缺乏覆盖检验，限制真实技术份额和未来绝对得分的解释。

因此本题的可执行建议为：在 A4 已观测配方及各问共同支持域内，先按目标权重与质量政策选择可行配方，再用 B1/B7 条件式和题面成本评估参数、数据与质量的预算分配；对超域方案、估算数据和跨来源能力换算显式返回限制状态。后续若取得同一模型、同一验证语料和相同 tokenization 下的 A/B/C 成对实验，可标定质量与 Loss 坐标并重新评价外推和技术贡献。
